\section{Post-Training}




\subsection{Multi-Teacher On-Policy Distillation (MOPD): A New Post-Training Paradigm}
\label{subsec:mopd-overview}

\begin{figure*}[t]
    \centering
    \includegraphics[width=\linewidth]{figures/post-training.pdf}
    \caption{Overview of MiMo-V2-Flash post-training stages.}
    \label{fig:post-training}
\end{figure*}


Modern language models increasingly rely on extensive post-training to enhance their intelligence and capabilities. However, current post-training pipelines face fundamental challenges: \textit{capability imbalance}, where improving one skill causes regressions in others (the ``see-saw'' effect), and \textit{learning inefficiency}, where existing approaches fail to fully leverage training signals when combining knowledge from multiple specialized models.

We propose Multi-Teacher On-Policy Distillation (MOPD), a unified post-training paradigm that addresses these challenges through a three-stage framework, as illustrated in Figure~\ref{fig:post-training}:

\paragraph{Stage 1: Supervised Fine-Tuning (SFT)} We establish foundational instruction-following capabilities through supervised learning on high-quality instruction-response pairs, enabling the model to understand and execute user requests across diverse domains.


\paragraph{Stage 2: Domain-Specialized Training} We train a suite of domain-specialized teacher models through independent RL optimization on focused tasks including agentic capabilities (search, coding, general tool use) and non-agentic tasks (mathematical reasoning, general reasoning, safety alignment). Each teacher achieves superior performance in its respective domain through targeted optimization with domain-specific reward signals.




\paragraph{Stage 3: Multi-Teacher On-Policy Distillation} Rather than merging model parameters or generating static offline datasets from experts, we formulate multi-teacher knowledge integration as an on-policy reinforcement learning process. The student model samples from its own evolving distribution and receives token-level supervision from domain-specific teachers through KL divergence rewards~\citep{minillm,Agarwal2023OnPolicyDO,lu2025onpolicydistillation}, effectively combining specialized capabilities without the traditional trade-offs~(Table~\ref{tab:mopd_eval}).

\begin{table}[htbp]
    \centering
    \small
    % \resizebox{0.65\textwidth}{!}{
    \begin{tabular}{l c c c c}
    \toprule
    \multirow{2}{*}{\textbf{Benchmark}}  & \textbf{Student} & \multirow{2}{*}{\textbf{Best Teacher}} & \textbf{Student} & \multirow{2}{*}{\textbf{$\Delta$(Student-Teacher)}} \\
                      & \textbf{Before MOPD} &  & \textbf{After MOPD} &  \\
    \midrule
AIME 2025                          &  89.3  &  93.9 \tiny{(RL)}  &  94.1 & $+$0.2 \\
HMMT Feb. 2025                          &  76.9  &  82.6 \tiny{(RL)}  &  84.4 & $+$1.8 \\
LiveCodeBench                &  77.5  &  82.6 \tiny{(RL)}  &  83.2 & $+$0.6 \\
MMLU-Pro                        &  84.7  &  84.7 \tiny{(Self)}  &  84.9 & $+$0.2 \\
GPQA-Diamond                    &  84.9  &  84.9 \tiny{(Self)}  &  84.3 & $-$0.6 \\
HLE \tiny{(w/o Tool)}                  &  21.2  &  21.2 \tiny{(Self)}  &  22.1 & $+$0.9 \\
Arena-Hard \tiny{(Hard Prompt)}        &  50.0  &  50.0 \tiny{(Self)}  &  54.1 & $+$4.1 \\
Arena-Hard \tiny{(Creative Writing)}   &  90.1  &  90.1 \tiny{(Self)}  &  86.2 & $-$3.9 \\
SWE-Bench Verified              &  67.8  &  74.2 \tiny{(RL)}  &  73.4 & $-$0.8 \\
Tau2-Bench                      &  75.9  &  79.6 \tiny{(RL)}  &  80.3 & $+$0.7 \\
Tau2-Bench \tiny{(Telecom)}              &  92.7  &  95.0 \tiny{(RL)}  &  95.3 & $+$0.3 \\
BrowseComp                      &  42.5  &  51.7 \tiny{(SFT)}  &  45.4 & $-$6.3 \\
    \bottomrule
    \end{tabular}
    % }
    \caption{Benchmark results of MOPD. The model types of best teachers are tagged, including RL, SFT, and the student model itself.}
    \label{tab:mopd_eval}
\end{table}






This unified framework offers several critical advantages over traditional post-training approaches:


\begin{itemize}
% [label=\textbf{(\roman*)},wide]
    \item \textit{Effective and Efficient.} 
    Unlike parameter merging or sequential training, which often trade off capabilities, MOPD preserves peak performance of the strongest teacher across all domains. Furthermore, on-policy distillation using dense, token-level rewards from teacher logits ensures stable credit assignment and rapid convergence. By learning from its own distribution, the student avoids the exposure bias and distribution mismatch common in off-policy methods trained on static datasets.
    \item \textit{Modular and Scalable.} The choice of teacher model is highly flexible: it can be a specialized RL-derived model with strong capabilities, a different SFT model, or even the student model itself. The decoupled design enables easy integration of new teachers without restructuring the entire pipeline. Moreover, the framework works seamlessly with existing outcome reward models (ORMs) and is especially advantageous for complex agentic tasks, where setting up independent training pipelines would otherwise be cumbersome.
    \item \textit{Iterative Co-Evolution.} MOPD naturally supports a teacher-student co-evolution cycle. Distilled student models can re-enter the specialized RL stage to produce stronger teachers, which in turn provide higher-quality supervision for the next generation of students, forming a self-reinforcing improvement cycle that enables sustained capability scaling.
\end{itemize}




In the following subsections, we detail each stage of the MOPD paradigm, beginning with supervised fine-tuning (\S\ref{subsec:post_train_sft}), followed by specialized RL for both agentic and non-agentic tasks~(\S\ref{subsec:scale_rl}), and conclude with the technical formulation of the multi-teacher distillation mechanism (\S\ref{subsec:mopd-technical}).


\subsection{Supervised Fine-Tuning (SFT)}
\label{subsec:post_train_sft}

The SFT stage serves as the foundation of our post-training pipeline, transforming the base model into a helpful assistant capable of following instructions and responding effectively across diverse tasks. This stage is crucial for activating the model's latent capabilities acquired during pre-training and aligning its outputs with desired formats and styles.

To achieve this, we curated millions of training samples spanning diverse domains, including general conversation, reasoning, coding, and agent tasks. These samples cover both thinking and non-thinking modes, with responses generated by our in-house domain-specialized model checkpoints. This diverse training mixture ensures comprehensive capability activation across the model's intended use cases.

Through preliminary experiments, we identified a critical stability metric for MoE SFT training: the number of parameters with zero gradients (num-zeros). This metric provides early warning signals for training instability: an increasing num-zeros indicates deteriorating load balance among experts, while a decreasing num-zeros suggests the model is significantly overfitting to the training data. Maintaining stable num-zeros throughout training is therefore essential for successful SFT. Furthermore, this stability is paramount for ensuring the robustness and convergence of the subsequent RL phase.

Our experiments reveal that num-zeros stability critically depends on two hyperparameters: the expert bias update rate and the $\epsilon$ parameter in the AdamW optimizer. Based on these findings, we configure our training with the following hyperparameters. We employ a cosine decay learning rate scheduler from $5.0 \times 10^{-5}$ to $5.0 \times 10^{-6}$, with a batch size of 128 and AdamW $\epsilon$ set to $1.0 \times 10^{-8}$. The MoE expert bias update rate is set to $1.0 \times 10^{-4}$, and the sequence auxiliary loss coefficient to $1.0 \times 10^{-6}$.
% 8 million SFT data
% lr: 5e-5 -> 5e-6 cosine decay, 5 epoch
% expert bias update: 1e-4
% sequence aux loss coeff: 1e-6
% adam epsilon: 1e-8


\subsection{Scaling Reinforcement Learning (RL)}
\label{subsec:scale_rl}

Reinforcement learning pushes model capabilities beyond what supervised fine-tuning alone can achieve. 
We employ different RL strategies depending on whether tasks involve agentic behavior, scaling both non-agentic and agentic RL training to maximize performance across diverse domains.

\subsubsection{Non-Agentic RL Training}
% RLVR + RLHF

Non-agentic RL training focuses on improving the model's performance on single-turn tasks, where the model generates a complete response without requiring interactive feedback or multi-step execution. The primary objective is to enhance the model's reasoning accuracy in verifiable domains (e.g., mathematics, coding, logic) while simultaneously aligning its outputs for helpfulness and safety in open-ended conversations.

Our approach to generating reward signals varies based on task characteristics. For domains with verifiable outcomes, we employ a hybrid verification system combining programmatic tools with an LLM judge to automatically assess correctness against curated problem-solution pairs. 
For subjective qualities such as helpfulness and safety, we implement a rubric-based framework where an advanced LLM judge evaluates responses against detailed rubrics and reference answers, producing granular reward signals that guide the model toward desired behaviors.

\subsubsection{Agentic RL Training}

\begin{table}[htbp]
\small
    \centering
    \begin{tabular}{lccc}
        \toprule
        \textbf{Agent Type} & \textbf{Number of Tasks} & \textbf{Environment} & \textbf{Prompt Source} \\
        \midrule
        Code Agent & $90\textrm{K}$ & Real & Real  \\
        Code Agent & $30\textrm{K}$ & Real & Synthesized \\
        Search Agent & $150\textrm{K}$ & Real & Synthesized \\
        General Agent & $50\textrm{K}$ & Synthesized & Synthesized \\
        \bottomrule
    \end{tabular}
    \caption{A summary of our training data composition across different agent types. We leverage both real-world and synthetically generated data to create a diverse set of tasks for training agents in various environments.}
    \label{tab:agent_tasks}
    
\end{table}
While non-agentic RL focuses on single-turn reasoning and generation, agentic RL trains the model to operate in interactive, multi-turn environments requiring planning, action execution, and adaptation based on feedback. We scale agentic RL along two critical dimensions: environment diversity and compute resources.

\paragraph{Scaling Agentic Environment Diversity} We construct a diverse suite of agentic training environments spanning code debugging, terminal operations, web development, and general tool use (Table~\ref{tab:agent_tasks}). Each environment targets distinct capabilities while sharing the common requirement of multi-step reasoning and execution. Below, we elaborate on the details for agentic environments.


\paragraph{Code Agent} We train on large-scale code agentic tasks derived from real-world GitHub issues, where the model operates in an agentic loop to read and edit files, execute commands, and receive rewards based on verifiable unit tests. Our key insight is that continuously scaling available tasks drives sustained improvements in code intelligence. To enable efficient RL training on over 100,000 code tasks, we develop two infrastructure components. First, we build an automated environment setup pipeline that provisions development environments from repository snapshots and packages them into containerized images, achieving 70\% success rate across 8 programming languages and supported by a large-scale Kubernetes cluster running over 10,000 concurrent pods. Second, we implement a lightweight agent scaffold that integrates seamlessly with Kubernetes, Docker, or local backends, exposing three atomic tools (\texttt{bash}, \texttt{str\_replace}, \texttt{finish}) that interact with execution backends solely via shell commands. This design eliminates server-based tool implementations and employs a minimal system prompt without predefined workflows, allowing the model to discover best practices during training.

\paragraph{Terminal Agent} Beyond GitHub issues, we strengthen terminal-based problem-solving capabilities using tasks sourced from Stack Overflow and Stack Exchange. We select materials requiring advanced technical expertise and transform them into computational tasks with corresponding queries, Dockerfiles, and test cases. After verifying environment installation and filtering by difficulty and reliability, we obtain approximately 30,000 queries with validated execution environments. Additional filtering based on pass rates removes tasks with unreliable correctness judgments or insufficient complexity for effective RL training.

\paragraph{Web Development Agent} To improve web development code generation, we build a real-world-grounded synthetic dataset paired with a multimodal verifier. We collect high-quality user-written web pages, execute generated code using Playwright to obtain rendered videos, and apply a multimodal visual discriminator to retain only high-quality samples, where video-based evaluation reduces visual hallucination compared to static screenshots. We reverse-engineer user queries from curated pages as seed prompts to synthesize large-scale RL data covering eight web categories that closely match real-world usage. Our vision-based verifier scores rollout executions from recorded videos, jointly evaluating visual quality, functional correctness, and executability to ensure rewards reflect both appearance and behavior.

\paragraph{General Agent} We develop two general agentic capabilities. Our search agent adopts a scaffold providing three core tools (search, open, find) for autonomous web exploration. We construct queries through recursive fact-graph expansion from seed entities, where difficulty scales with relation chain depth and detail obfuscation, enabling automated generation of challenging search problems with verifiable answers. Our function-calling agent trains on synthetic application environments with custom toolsets constructed by generating tool-call graphs based on explicit data dependencies (direct input-output relationships) and implicit logical dependencies (reasoning about hidden system states), requiring both data propagation and state inference capabilities.


\paragraph{Scaling Agentic Compute}
Training on the previous diverse set of agentic environments (Table~\ref{tab:agent_tasks}), we find that scaling agentic RL compute not only boosts code-agentic performance but also generalizes effectively to other task types. Figure~\ref{fig:swebench_scaling} shows the RL training curve for our code-agent, where the model performed on-policy rollouts and updates across approximately 120K environments. 
This scaling significantly improves upon the SFT base model's performance on SWE-Bench-Verified and SWE-Bench-Multilingual. Moreover, Figure~\ref{fig:swe_generalize} demonstrates that large-scale code-agentic RL training generalizes effectively to other agentic tasks, as well as math, code, and general reasoning benchmarks, suggesting that agentic training develops broadly transferable problem-solving capabilities.

\begin{figure*}[t!]
    \centering
    \includegraphics[width=0.95\linewidth]{figures/swebench_scaling.pdf}
    \caption{Code-agentic RL scaling curves. The X-axis represents total interactive environments consumed during on-policy RL rollouts; the Y-axis shows resolved rates on SWE-Bench-Verified and SWE-Bench-Multilingual.}
    \label{fig:swebench_scaling}
\end{figure*}

\begin{figure*}[t!]
    \centering
    \includegraphics[width=1.0\linewidth]{figures/swe_generalization.pdf}
    \caption{Generalization of code-agentic RL training to other task domains.}
    \label{fig:swe_generalize}
\end{figure*}


\subsection{Technical Formulation of MOPD}
\label{subsec:mopd-technical}

Having established the foundation through SFT and trained specialized teachers through domain-specific RL, we now formalize the multi-teacher on-policy distillation mechanism that integrates these specialized capabilities into a unified student model.

Specifically, we cast multi-teacher distillation as an on-policy reinforcement learning objective. Let $\pi_\theta$ denote the target student policy optimized in the training engine, $\mu_\theta$ denote the student sampling policy adopted in the inference engine, and $\pi_{\text{domain}_x}$ denote the teacher policy specialized for the domain of prompt $x$ sampled from distribution $\mathcal{D}$. Let $\text{sg}[\cdot]$ denote the stop-gradient operator. The reverse KL divergence loss between student and teacher is defined as:
\begin{equation}
    \mathcal{L}_\text{reverse-KL}(\theta)
  =-\mathbb{E}_{x\sim\mathcal{D},  y_{t}\sim \pi_\theta(\cdot|x, y_{<t})}\log\frac{ \pi_{\text{domain}_x}(y_{t}|x, y_{<t})}{ \pi_\theta(y_{t}|x, y_{<t})}.
\end{equation}
The gradient is:
\begin{equation}
    \nabla_{\theta}\mathcal{L}_\text{reverse-KL}(\theta)
  =-\mathbb{E}_{x\sim\mathcal{D},  y_{t}\sim \pi_\theta(\cdot|x, y_{<t})}\left[\log\frac{ \pi_{\text{domain}_x}(y_{t}|x, y_{<t})}{ \pi_\theta(y_{t}|x, y_{<t})} \nabla_\theta \log \pi_\theta(y_{t}|x, y_{<t})\right].
\end{equation}
Following~\citet{IcePop2025}, we apply training-inference importance sampling and discard tokens that exhibit large discrepancies. We then define the surrogate loss of MOPD as:
\begin{equation}
\mathcal{L}_{\text{MOPD}}(\theta) = -\mathbb{E}_{x \sim\mathcal{D}, y \sim \mu_{\theta}(\cdot|x)}
\left[\frac{1}{|y|}\sum_{t=1}^{|y|}
  w_{t} \hat{A}_{\text{MOPD},t}\log\pi_{\theta}(y_{t}|x, y_{<t}) \right],
\end{equation}
where
\begin{equation}
    w_{t}(\theta)=
\begin{cases}
\text{sg}\left[\frac{\pi_\theta(y_{t}|x, y_{<t})}{\mu_\theta(y_{t}|x, y_{<t})}\right], & \epsilon_{\text{low}} \le \frac{\pi_\theta(y_{t}|x, y_{<t})}{\mu_\theta(y_{t}|x, y_{<t})} \le \epsilon_{\text{high}}, \\
0,  & \text{other},
\end{cases}
\quad
    \ \hat{A}_{\text{MOPD},t} = \text{sg}\left[\log\frac{ \pi_{\text{domain}_x}(y_{t}|x, y_{<t})}{ \pi_\theta(y_{t}|x, y_{<t})}\right].
\label{eq:advantage_calculation}
\end{equation}

By default, we combine the advantages of MOPD with other types of advantages, such as those computed using Outcome Reward Models (ORMs), including GRPO~\citep{shao2024deepseekmath}. 
Let $\hat{A}_{\text{ORM}}$ denote the advantages computed by the ORMs; the final advantages are given by:
\begin{equation}
    \hat{A}_{\text{MOPD},t} = \text{sg}\left[\log\frac{ \pi_{\text{domain}_x}(y_{t}|x, y_{<t})}{ \pi_\theta(y_{t}|x, y_{<t})}\right] + \alpha \hat{A}_{\text{ORM}}.
\end{equation}





\begin{figure}[t]
    \centering
    \includegraphics[width=0.48\linewidth]{figures/post-train-compare-AIME25.pdf}
    \includegraphics[width=0.48\linewidth]{figures/post-train-compare-LiveCodeBench.pdf}
    \caption{Comparison of different post-training methods on math and code tasks. Three lines represent training RL with ORM, MOPD without outcome rewards (MOPD w/o ORM), and MOPD. }
    \label{fig:compare-post-train}
\end{figure}

Figure~\ref{fig:compare-post-train} demonstrates the effectiveness of MOPD compared to traditional post-training approaches. On mathematical reasoning (AIME 2025) and coding (LiveCodeBench) benchmarks, MOPD successfully preserves and combines specialized capabilities from multiple teachers, achieving performance that matches or exceeds the strongest teacher in most domains.


\subsection{Evaluations}
\label{subsec:post_train_eval}


\subsubsection{Evaluation Setup}

We evaluate \mimoflash{} on MMLU-Pro~\citep{wang2024mmlu}, GPQA-Diamond~\citep{rein2024gpqa}, HLE Text-only~\citep{phan2025humanity}, AIME 2025~\citep{AIME}, LiveCodeBench (2024.08-2025.04)~\citep{jain2024livecodebench},  HMMT Feb. 2025~\citep{balunovic2025matharena}, Arena-Hard~\citep{arenahard2024}, LongBench V2~\citep{bai2025longbench}, MRCR~\citep{vodrahalli2024michelangelo} (\{2,4,8\}-needles, maximum 128K), SWE-Bench Verified~\citep{jimenez2024swebench}, SWE-Bench Multilingual~\citep{yang2025swesmith}, Terminal-Bench, BrowseComp~\citep{wei2025browsecomp}, $\tau^2$-Bench~\citep{barres2025tau}.

\subsubsection{Evaluation Results}

We illustrate the evalution results in Table~\ref{tab:main_eval}.
\mimoflash{} achieves performance comparable to that of Kimi-K2-Thinking and DeepSeek-V3.2-Thinking on most reasoning benchmarks.
The model also maintains competitive general writing capabilities, enabling it to generate high-quality responses on open-ended tasks.
In long context evaluations, our model surpasses Kimi-K2-Thinking, a significantly larger full global attention LLM, highlighting the strong long-context capabilities of our hybrid SWA architecture.

Notably, \mimoflash{} achieves 73.4\% on SWE-Bench Verified, outperforming all open-source competitors and approaching the performance of GPT-5-High.
On SWE-Bench Multilingual, our model resolves 71.7\% issues, establishing it as the most capable open-source LLM for software engineering tasks.
These results underscore the effectiveness of our ultra-scaled agentic RL training. On Terminal Bench, the model also delivers a competitive score. 

In search agent evaluation, \mimoflash{} scores 45.4 on BrowseComp, and is further boosted to 58.3 with the context management method outlined in Appendix~\ref{sec:context_management}.
For general tool-use on $\tau^2$-Bench, we employ DeepSeek-V3.2 as the user agent, achieving category scores of 95.3 (Telecom), 79.5 (Retail), 66.0 (Airline).

Taken together, these results validate the effectiveness of our ultra-large-scale RL training within the MOPD post-training paradigm, and highlight the models strong potential for real-world coding, reasoning, and agentic workflows.


\begin{table}[htbp]
    \centering
    \small
    \resizebox{0.99\textwidth}{!}{
    % \setlength{\tabcolsep}{4.5pt}
    \begin{tabular}{l | c | c c c c c}
    \toprule
    \multirow{2}{*}{\textbf{Benchmark}} & \textbf{MiMo-V2} & \textbf{Kimi-K2} & \textbf{\textbf{DeepSeek-V3.2}} & \textbf{Gemini-3.0} & \textbf{Claude} & \textbf{GPT-5} \\
    & \textbf{Flash} & \textbf{Thinking} & \textbf{Thinking} & \textbf{Pro} & \textbf{Sonnet 4.5} & \textbf{High} \\
    \midrule
    \multicolumn{6}{l}{\textbf{Reasoning}} \\
    MMLU-Pro & 84.9 & 84.6 & 85.0 & 90.1 & 88.2 & 87.5 \\
    GPQA-Diamond & 84.3 & 84.5 & 82.4 & 91.9 & 83.4 & 85.7 \\
    HLE {\tiny (no tools)}  & 22.1 & 23.9 & 25.1 & 37.5 & 13.7 & 26.3 \\
    AIME 2025 & 94.1 & 94.5 & 93.1 & 95.0 & 87.0 & 94.6 \\
    HMMT Feb. 2025 & 84.4 & 89.4 & 92.5 & 97.5 & 79.2 & 88.3 \\
    LiveCodeBench-v6 & 85.1 & 83.1 & 83.3 & 90.7 & 64.0 & 84.5 \\
    \midrule
    \multicolumn{6}{l}{\textbf{General Writing}} \\
    Arena-Hard {\tiny (Hard Prompt)} & 54.1 & 71.9 & 53.4 & 72.6 & 63.3 & 71.9 \\
    Arena-Hard {\tiny (Creative Writing)} & 86.2 & 80.1 & 88.8 & 93.6 & 76.7 & 92.2 \\
    % Creative Writing &  & 82.4 &  & 81.5 & 80.7 & 84.0 \\
    \midrule
    \multicolumn{6}{l}{\textbf{Long Context}} \\
    LongBench V2 & 60.6 & 48.1 & 58.4 & 65.6 & 61.8 & - \\
    MRCR & 45.7 & 44.2 & 55.5 & 89.7 & 55.4 & - \\
    \midrule
    \multicolumn{6}{l}{\textbf{Code Agent}} \\
    SWE-Bench Verified & 73.4 & 71.3 & 73.1 & 76.2 & 77.2 & 74.9 \\
    SWE-Bench Multilingual & 71.7 & 61.1 & 70.2 & - & 68.0 & 55.3 \\
    Terminal-Bench Hard &30.5  &30.6 & 35.4 & 39.0 & 33.3 & 30.5 \\
    Terminal Bench 2.0 &38.5  & 35.7 & 46.4 & 54.2 & 42.8 & 35.2 \\
    \midrule
    \multicolumn{6}{l}{\textbf{General Agent}} \\
    BrowseComp & 45.4 & - & 51.4 & - & 24.1 & 54.9 \\
    BrowseComp {\tiny (w/ Context Manage)} & 58.3 & 60.2 & 67.6 & 59.2 & - & - \\
    $\tau^2$-Bench & 80.3 & 74.3 & 80.3 & 85.4 & 84.7 & 80.2  \\
    \bottomrule
    \end{tabular}
        }
    \caption{
        Comparison between \mimoflash{} and open/closed models.
    }
    \label{tab:main_eval}
\end{table}










\subsection{RL Infrastructures}

Our RL (and MOPD) infrastructure uses SGLang~\citep{zheng2024sglang} as the inference engine and Megatron-LM~\citep{shoeybi2019megatron} as the training engine.
We adopt FP8 for both training and inference.
To enable stable, efficient, and flexible RL training, we implement three extended modules: Rollout Routing Replay~\citep{ma2025stabilizing} (Sec~\ref{sec:posttrain:rlinfra:r3}), Data Scheduler (Sec~\ref{sec:posttrain:rlinfra:data}), and Tool Manager combined with Toolbox (Sec~\ref{sec:posttrain:rlinfra:tool}).

\subsubsection{Stablized Training via Rollout Routing Replay (R3)}
\label{sec:posttrain:rlinfra:r3}

MoE models suffer from inconsistent expert routing across rollout and training due to numerical precision issues~\citep{he2025nondeterminism,yao2025offpolicy}.
%
We propose Rollout Routing Replay (R3)~\citep{ma2025stabilizing} to train RL using the same routed experts from rollout, making its overhead negligible through optimized data types and communication overlapping.
%
For multi-turn agent training, we employ a request-level prefix cache during rollout.
%
This cache stores KVCache and MoE routed experts from prior turns, allowing them to be reused for subsequent generation steps of the same request.
%
Unlike the commonly-used radix cache in current inference engines, our request-level prefix cache avoids re-prefilling or inter-request output cache sharing, ensuring sampling consistency for routed experts.



\subsubsection{Data Scheduler}
\label{sec:posttrain:rlinfra:data}

For MiMo-V2-Flash, we extend the Seamless Rollout Engine~\citep{xia2025mimo} and implement a Data Scheduler to seamlessly schedule fine-grained sequences instead of micro-batches, addressing GPU idleness in distributed MoE training.
%
In dynamic sampling, as sequences return for reward computation, we reference historical pass rates and, if necessary, assign new prompts to GPUs with load balancing.
%
We integrate partial rollout~\citep{team2025kimi,fu2025areal} to partition overlong trajectories across steps, while limiting staleness and the proportion of partial samples in each batch.
%
By employing staleness-aware truncated importance sampling for partial rollout, we significantly accelerate RL training without sacrificing model quality.

The Data Scheduler supports data source-specific configurations (sample quota, scheduling priority, length limits, temperature) and fits pass rates to accept samples by configured ratios.
%
Priority-based scheduling overlaps reward computation and inference across data sources with different time patterns, ensuring high GPU utilization.

\subsubsection{Toolbox and Tool Manager}
\label{sec:posttrain:rlinfra:tool}

We implement Toolbox and Tool Manager to tackle global resource contention and local inefficiency in RL agent training.
%
These modules leverage Ray~\citep{moritz2018ray} for efficient scheduling.
%
Toolbox acts as the centralized resource allocator, enforcing resource quota and QPS limits for tools across concurrent tasks.
%
It adopts fault-tolerant Ray actor pools, which eliminate cold-start delays.
%
Integrated with the rollout engine, Tool Manager coordinates with Toolbox to accelerate training through environment pre-warming and sequence-level asynchronous reward computation.
%
It maintains training stability through timeout recovery and real-time monitoring.
%
By disaggregating the tool management and rollout workflow, Toolbox isolates task-specific logic from system-wide policies, enabling modular extensibility without compromising stability.



